\chapter[Introdução]{Introdução}
\label{cap:introducao}

O registro eletrônico em saúde concentra grande parte da informação clínica
sob a forma de texto livre. Notas de evolução, sumários de alta e laudos são
redigidos em linguagem natural, com abreviações idiossincráticas, ausência
frequente de pontuação padronizada e construções negativas características do
discurso médico. Essa riqueza textual, embora valiosa para pesquisa e apoio à
decisão, permanece majoritariamente inacessível a consultas estruturadas, o que
motiva o desenvolvimento de técnicas de Processamento de Linguagem Natural (PLN)
capazes de transformar texto clínico em conhecimento estruturado.

No contexto da língua portuguesa, o avanço dessas técnicas foi historicamente
limitado pela escassez de recursos anotados. A disponibilização do corpus
\textbf{SemClinBr} \cite{oliveira_2022_semclinbr}, primeiro corpus clínico
anotado semanticamente para o português, com mil documentos provenientes de
múltiplas instituições e especialidades, criou condições para investigar tarefas
antes restritas ao inglês, entre elas a extração de relações entre entidades
clínicas.

\section{Problema de pesquisa}

A extração de entidades nomeadas identifica menções a conceitos clínicos
(sintomas, doenças, achados, procedimentos), mas isoladamente não captura como
esses conceitos se articulam no enunciado. Saber que um documento menciona
``febre'' não é suficiente: é preciso distinguir se a febre está \emph{associada}
ao quadro do paciente ou se está sendo \emph{negada} (``paciente afebril'').
Esse vínculo é o objeto da \textbf{extração de relações} (\foreignlanguage{english}{Relation Extraction}, RE).

Este trabalho investiga a seguinte questão: \emph{em que medida é possível
classificar automaticamente as relações entre entidades clínicas anotadas no
SemClinBr, dentro do espaço de rótulos \texttt{negation\_of},
\texttt{associated\_with} e \texttt{no\_relation}, e qual o ganho de modelos
aprendidos sobre abordagens léxicas?} O desafio é agravado por dois fatores
inerentes ao corpus: o forte desbalanceamento entre as classes de relação e o
fato de o SemClinBr anotar apenas relações positivas, exigindo a geração
controlada de exemplos negativos.

\section{Motivação}

A negação clínica é um fenômeno linguístico de alto impacto: tratar um achado
negado como presente pode inverter completamente a interpretação de um
prontuário. Em português, a negação é fortemente lexicalizada (``não'', ``sem'',
``nega'', ``ausência de''), o que sugere que mesmo abordagens simples podem
estabelecer um patamar de desempenho relevante e, ao mesmo tempo, evidenciar as
lacunas que justificam modelos contextuais. Construir esse percurso de forma
reprodutível — desde o texto bruto até as métricas finais — é, em si, uma
contribuição metodológica, dado que a falta de reprodutibilidade é uma das
principais críticas dirigidas à literatura de PLN biomédico.

\section{Objetivos}

\subsection{Objetivo geral}

Desenvolver, de forma reprodutível, um \emph{pipeline} determinístico de extração de
relações entre entidades clínicas em português sobre o corpus SemClinBr ---
considerando o espaço de rótulos \texttt{negation\_of}, \texttt{associated\_with} e
\texttt{no\_relation} --- e avaliar, sob protocolo estatístico rigoroso (macro-F1 e
F1 por classe com intervalos de confiança por \foreignlanguage{english}{bootstrap} e
teste de significância pareado), uma cadeia de três \emph{baselines} de complexidade
crescente: classe majoritária (B0), regras léxicas (B1) e regressão logística com
características (B2). O ajuste fino de um modelo contextual de domínio clínico
(BioBERTpt) é posicionado como trabalho futuro imediato, para o qual a cadeia
estabelecida fornece patamares de comparação.

\subsection{Objetivos específicos}

\begin{itemize}
  \item Construir um parser fiel ao formato XML do SemClinBr que produza uma
        representação canônica determinística do corpus.
  \item Caracterizar o corpus por meio de análise exploratória, quantificando a
        distribuição das relações e a distância entre entidades relacionadas.
  \item Definir uma partição treino/desenvolvimento/teste em nível de documento,
        estratificada e congelada, de modo a evitar vazamento de dados.
  \item Implementar um gerador determinístico de pares candidatos compartilhado por
        todos os modelos, garantindo a comparabilidade dos resultados.
  \item Estabelecer e avaliar três \emph{baselines} --- classe majoritária (B0),
        regras léxicas inspiradas no algoritmo NegEx \cite{chapman_2001_negex} (B1)
        e regressão logística com características léxicas e estruturais (B2) ---,
        reportando macro-F1 e F1 por classe com intervalos de confiança por
        \foreignlanguage{english}{bootstrap} e comparação por teste de McNemar.
  \item Preparar a infraestrutura para o ajuste fino do BioBERTpt como modelo
        proposto, a ser executado como continuidade imediata do trabalho.
\end{itemize}

\section{Organização do trabalho}

O Capítulo~\ref{cap:referencial} apresenta os fundamentos de PLN clínico,
extração de relações e modelos baseados em \foreignlanguage{english}{transformers}.
O Capítulo~\ref{cap:relacionados} discute trabalhos relacionados, com ênfase em
recursos para o português. O Capítulo~\ref{cap:metodologia} descreve o corpus, o
pipeline de processamento, a geração de pares candidatos e o protocolo de
avaliação. O Capítulo~\ref{cap:prototipo} detalha a proposta de modelo. O
Capítulo~\ref{cap:experimentos} reporta os experimentos e resultados. O
Capítulo~\ref{cap:discussao} analisa limitações e ameaças à validade, e o
Capítulo~\ref{cap:conclusao} sintetiza as contribuições e aponta trabalhos
futuros.
